
Throughout this section, we will explain the followed methodology for these probes. Firstly, we will not be doing real space analysis, instead, in this work we will be working in the harmonic space by using power spectrum analysis. More specifically, our purpose is to study the Angular Power Spectrum (APS) of the synchrotron emission.

\subsection{Power spectra analysis}
\label{sec:power_spectrum}

The APS is a mathematical representation of the distribution of power as a function of angular scale, allowing us to study energy fluctuations 
at different angular scales. This is a very powerful tool for understanding statistical properties by extracting information about the underlying 
physical processes giving rise to these fluctuations.
\paragraph*{}

The power spectrum is computed as a function of $a_{\ell m}$, the spherical harmonic coefficients. Noting that these coefficients are a function of the multipole $\ell$, which represents the different angular scales: large angular scales correspond to low $\ell$ values, while small angular scales correspond to high $\ell$ values. For example, $\ell = 200$ represents approximately $1^{\circ}$ scales, which is the angular resolution for QUIJOTE. The mathematical representation of the APS is given by:
\paragraph*{}


\begin{equation}
    \label{eq:cl_auto}
    C_{\ell} = \left < \frac{\sum_{m}^{} a_{\ell m} (a_{\ell m})^{*} }{2\ell +1} \right > = \left < \frac{\sum_{m}^{} |a_{\ell m}|^2}{2\ell +1} \right >.
\end{equation}

This is only for the computation of auto-power spectrum, which implies only one frequency map at time. For auto-spectra, we are taking into account 
only the three maps (I, Q and U) for each frequency band separately. But in this case, we will also incorporate the cross-power spectrum data, which give us in a simplified way, the common signal between two different frequency bands. This is a useful data addition in order to improve our results. To obtain 
these cross-spectra, we must compute the $C_{\ell}$ on a different way. Now, two different $a_{\ell m}$ coefficient will be implied. Therefore, the mathematical 
representation for this cross-spectrum is:

\begin{equation}
    \label{eq:cl_cross}
    C_{\ell}^{X - Y} = \left < \frac{\sum_{m}^{} a_{\ell m}^{X} (a_{\ell m}^{Y})^{*} }{2\ell +1} \right >,
\end{equation}

where X and Y represent the two frequency bands that we are correlating. Note that if X and Y are the same frequency band, this expresion simplifies as the 
auto-spectrum analysis. The power spectrum computation will be carried out through the scientific \python\ library, \namaster\footnote{\url{https://namaster.readthedocs.io/en/latest/}} (\citealp{Namaster}).
\paragraph*{}

\subsection{Beam, noise and colour corrections}
\label{sec:corrections}

The $C_{\ell}$ are computed using real sky data, containing noise and other contamination sources. In the ideal case we would obtain $C_{\ell} \ \propto \ |a_{\ell m}^{theo}|^2$, but these coefficients must be corrected due to different effects:

\subsubsection{Window function}

Window functions appear as a multiplicative function to the theoretical $a_{\ell m}^{theo}$ coefficients. Throughout this work, we are correcting the pixel and beam window functions. The first one appears as a consequence of \healpix\ map pixelization. When the data is converted to a map, the $a_{\ell m}$ amplitude decreases as a function $w_{\ell m}$, this function increases as we study smaller scales or as $\ell$ increases. This function only depends on the map $N_{side}$ resolution. We are using $N_{side} = 512$ maps, so this correction term is not very significant. 
\paragraph*{}

In contrast, the beam window function, denoted as $B_{\ell}$, is a consequence of the instrumental window function, depending on the instrument resolution. It is a multiplicative factor to the $a_{\ell m}$ coefficients, this term grows exponentially at the resolution scales of the instrument, being very dominant at small angular scales. So the actual dependence of the auto-power spectrum becomes $C_{\ell} \ \propto \ |a_{\ell m}^{theo}|^2 |w_{\ell m}|^2 B_{\ell}^2$, however, for the cross-spectrum we obtain: $C_{\ell}^{X-Y} \ \propto \ a_{\ell m}^{X,theo} (a_{\ell m}^{Y,theo})^{*} |w_{\ell m}|^2 B_{\ell}^{X} (B_{\ell}^{Y})^{*}$. As the two maps implied must have the same pixelization resolution, both pixel window functions are the same, therefore, we can simplify this correction term.

\subsubsection{Noise power spectrum}

The real sky maps contain an instrumental noise component that must be corrected from the data. In order to obtain this component, we will 
follow the procedure explained in \citealp{QUIJOTE_IV}. We obtain a Half Mission Difference Map (HMDM) for each QUIJOTE and Planck frequency band used for this 
probe, produced for each case, as:

\begin{equation}
    \label{eq:hmdm}
    n = \frac{h_1 - h_2}{\sqrt{ (w_1 + w_2) (w_1^{-1} + w_2^{-1}) }}. 
\end{equation}

Here $w_{1}$ and $w_{2}$ are the individual weight maps of the null tests half 1 ($h_{1}$) and half 2 ($h_{2}$), respectively. They are computed as 
$w_{i} = 1/\sigma_{i}$, with $i = 1, 2$. Defined in this way, Equation \ref{eq:hmdm} provides a map with similar noise levels as the residual 
noise for the weighted-sum of the two halves (see e.g. \citealp{Planck_2014b}, \citealp{Planck_2016f}). However, for WMAP data we will obtain our 
noise power spectrum as the mean spectrum between 100 power spectra obtained from noise simulations. These simulations are explained 
in Section \ref{sec:noise_simulations}.
\paragraph*{}

Now, by computing and correcting the window functions effects, we can obtain the noise power spectrum for the auto-spectrum. The estimated $C_{\ell}$ 
on sky is therefore: 

\begin{equation}
    \label{eq:corrected_cl}
    C_{\ell} = (C_{\ell}^{sky} - N_{\ell}) /( |w_{\ell m}|^2 B_{\ell}^2),
\end{equation}

where $N_{\ell}$ is the HMDM auto-power spectrum in the case of QUIJOTE and Planck data, or the mean power spectrum obtained from noise simulations in the case 
of WMAP data. For the cross-power spectrum no noise correction will be needed. As the different frequency bands that 
we will be cross-correlating have been measured by independent instruments, their noise fluctuations are not correlated. Taking into account this correction only for 
auto-spectrum analysis. Therefore, cross-spectra corrections are applied as:

\begin{equation}
    \label{eq:corrected_cross}
    C_{\ell}^{X-Y, corr} = C_{\ell}^{X-Y} / ( |w_{\ell m}|^2 B_{\ell}^{X} B_{\ell}^{Y}).
\end{equation}

Here $X$ and $Y$ are referring the two different frequency bands that we are cross-correlating. 

\subsubsection{Colour corrections}
\label{sec:colour_corrections}

Colour corrections are mainly associated with the proper characterization of the instrument bandpass response. In Sections \ref{sec:full_sky_analysis} 
and \ref{sec:local_sky_analysis}, we will be performing Markov chain Monte Carlo (MCMC) fittings. In order to account for the colour corrections effect, 
we will be following
the same procedure as explained in \citealp{Planck_2020_cc}. Obtaining our colour correction terms as an iterative process carried out by fitting our data 
without taking into account this correction on a first step, and then, using the obtained spectral indices, we obtain the colour corrections terms for 
each fitting through the \fastcc \footnote{\url{https://github.com/mpeel/fastcc}} \python\ and \IDL\ code (\citealp{Fastcc}). At this point, we can now 
repeat this procedure taking into account colour correction in the MCMC fitting process.

\subsection{Masks}
\label{sec:masks}

In this work we will not be doing full sky analysis, instead we will be presenting Northern sky and local region analyses. For this purpose, through 
\python\ libraries \healpix\ and \namaster, we construct a set of customized masks focusing on different sky regions. The \namaster\ power spectrum computation 
requires that all masks must have smooth and derivable selection borders (\citealp{Namaster}), this is done by following a mask apodization process. In this 
work we will be using the cosine function apodization (C2 apodization type) with a 5$^{\circ}$ smoothing border as implemented in the \namaster\ library. 
The Northern sky regions will be presented in Section \ref{sec:full_sky_analysis}, while local regions are discussed in Section \ref{sec:local_sky_analysis}.

\subsection{Simulations}
\label{sec:simulations}

From Equation \ref{eq:corrected_cl}, it is possible to deduce that the errors of the power spectrum obtained for each polarization mode and frequency band are computed as the quadratic sum of the error inherent to the sky signal map and the error of the noise map. These are estimated as the standard deviation between $N$ realistic sky simulated maps power spectra, including noise; and the standard deviation between the simulated noise power spectra. In this work, we choose $N=100$. In this way, we can quantify the errors inherent to the real sky power spectrum and to the subtracted noise power spectrum component. So, in order to estimate the errors of our power spectra analyses, we need a set of simulations, both for a realistic sky emission model and for the instrumental noise in intensity and polarization. 


\subsubsection{Sky signal}
\label{sec:sky_simulations}

Python Sky Model package\footnote{\url{https://pysm3.readthedocs.io/en/latest/}} (\pysm; \citealp{Pysm}) is a \python\ library used for the sky modelling of many CMB analysis (e.g. \citealp{Planck_2020_cc}, \citealp{Martire_2022}). This library enables us to achieve a realistic simulated sky emission map as seen by 
each instrument that we are interested in, but without a noise component. In our case, four sky simulations are needed, centered at 11.1- and 12.9-GHz for the horn 3 QUIJOTE-MFI, at 22.8 GHz for the WMAP K-band, and the last one at 28.4 GHz for Planck 30 GHz channel. For each one of these simulations we are creating a map with the same angular resolution as each instrument. These resolutions are given in the previous Section \ref{sec:data}.
\paragraph*{}

The selected sky model includes all known cosmic foregrounds. These are for the polarized components: synchrotron emission, described as a power law 
scaling with non spatially varying spectral index (\textbf{s1}), thermal dust, modelled as a single-component modified black body (\textbf{d1}), and a 
lensed CMB realization, computed using Taylens\footnote{\url{https://github.com/amaurea/taylens}}, a code to compute a lensed CMB realization 
using nearest-neighbour Taylor interpolation (\citealp{Taylens}). This code takes, as an input, a set of unlensed $C_{\ell}$ generated using 
CAMB\footnote{\url{https://www.camb.info/}} (\textbf{c1}). Also, in order to obtain the most realistic sky map as possible, we decided to include intensity 
radiations foregrounds, even if in this case this is not necessary as we are carrying out a polarization probe. These foregounds are: AME modeled as 
a sum of two spinning dust populations based on the \commander \footnote{\url{https://commander.bitbucket.io/}} code (\citealp{Planck_2015_commander}) 
(\textbf{a1}), and free-free emission, modelled using the analytic model assumed in the Commander fit to the Planck 2015 data (\textbf{f1}). 


\subsubsection{Simulated noise maps}
\label{sec:noise_simulations}

Noise simulations are specific for each instrument. Therefore, these simulations must be obtained from completely independent processes for each 
frequency band. In the case of QUIJOTE 11 and 13 GHz channels, as these two bands are measured using the same horn of the MFI, their noise fluctiations are not 
completely independent. In this case, we need our simulations to account for the measured anisotropic behaviour, spatial correlations, 
and the correlations between these two frequency channels ($\sim 60 - 80 \%$ in intensity, and $\sim 20 \%$ in polarization, \citealp{QUIJOTE_IV}). 
In this work we will be using the QUIJOTE noise simulations obtained as explained in \citealp{QUIJOTE_IV}. These simulations 
have been provided by the author. 
\paragraph*{}

The WMAP K-band noise simulations have been obtained following the methodology explained in \citealp{WMAP_maps}. Making 100 noise simulation maps through 
random Gaussian simulations as WMAP data is not affected by $1/f$ noise. These simulations must have zero mean value and $\vec{\sigma}$ deviation, 
different for each pixel. We are taking into account that this noise decreases by the square root of the number of observations of each individual pixel. 
So we can obtain the deviation for each pixel as:

\begin{equation}
    \label{eq:wmap_noise}
    \vec{\sigma} = \frac{\sigma_{0}}{\sqrt{N_{obs}}},
\end{equation}

where $\sigma_{0}$ represents the radiometer noise for an individual sample, and $N_{obs}$ is the effective number of observations for each pixel, 
and ignoring the covariance term from QU. Once we have obtained this deviation, we can compute 100 random Gaussian simulations for WMAP noise. 
As for this experiment, there are no half-mission maps, the noise power spectrum subtracted in Equation \ref{eq:corrected_cl} will be the mean of the power spectra for 100 noise simulations.
\paragraph*{}

Finally, for Planck 30 GHz channel, 100 instrumental noise simulations for each data release have been downloaded from Planck Legacy Archive (these simulations take up 47 GB of space for the PR3 data and 15 GB for the PR4 data). These noise simulations have been described in \citealp{Planck_PR4_maps} in the case of PR4 data and in \citealp{Planck_pr3_maps} for the PR3 maps. These noise maps have been made in order to provide a good estimation of instrumental noise by following the process shown in \citealp{Planck_2014b}, using an implementation 
of a Markov Chain Monte Carlo (MCMC) approach to estimate basic noise properties.
\paragraph*{}


